\chapter{The open-problem ledger}
\label{app:ledger}

The program's rule 5 requires every dive that discovers a new open problem to append it to the backlog with a
one-line rationale. Eighty-eight such items existed at compilation. \Cref{tab:ledger} lists them in numerical
order with the dive that opened each, its status on 2026-09-07 as recorded in the index (``resolved'',
``priced'' and ``discharged'' are the index's own words; ``partly'' marks items the index says were advanced
but not closed), and the chapter of this compendium in which the item is discussed. The one-line rationales
are in the chapters' ``Further exploration'' sections; the index file carries the full text.

Fourteen items were closed by later dives: \BL{1}, \BL{2}, \BL{73} and \BL{75} by the Phase~D identification
chain; \BL{5} and \BL{6} by \dive{04}; \BL{20} by \dive{08}; \BL{21} by \dive{09}; \BL{42} by \dive{14};
\BL{57} and \BL{58} by \dive{19} (\BL{56} partly); \BL{70} was priced rather than solved by \dive{21}; and
\BL{82} was discharged for dives 20--22 only. The index names three candidates for run 23: computing the
shape-prior ledger on a real fitted MMM's Hessian together with the flighting-calendar rank audit
(\BL{81} + \BL{74} + \BL{86}), the response-curve estimand for prior pricing (\BL{85}), and intermodulation
under synergy (\BL{79}).

\small
\begin{longtable}{@{}r p{3.35in} c p{1.55in} c@{}}
\caption{The open-problem ledger: every backlog item (BL) spawned by the program, in numerical order, with the dive that opened it, its status as of 2026-09-07, and the chapter of this compendium in which it is discussed.}\label{tab:ledger}\\
\toprule
BL & Problem & Opened by & Status & Chapter \\
\midrule
\endfirsthead
\toprule
BL & Problem & Opened by & Status & Chapter \\
\midrule
\endhead
\bottomrule
\endfoot
1 & Seasonality-orthogonal perturbation design & 01 & resolved by \dive{20} & 1 \\
2 & The price of identification & 01 & resolved by \dive{20} & 1 \\
3 & Calibration under model misspecification & 02 & open & 2 \\
4 & A specification test for transport & 02 & open & 2 \\
5 & The EVSI power plateau --- boundaries & 03 & resolved by \dive{04} & 3 \\
6 & Second-order EVSI scoring & 03 & resolved by \dive{04} & 3 \\
7 & Power--cost EVSI frontier & 03 & partly advanced by \dive{04} & 3 \\
8 & Optimal switch-experiment sizing & 04 & open & 3 \\
9 & The k-experiment portfolio law & 04 & open & 3 \\
10 & LQC robustness under misspecification & 04 & open & 3 \\
11 & Belief-independent risk ledgers & 05 & open & 4 \\
12 & The low-drift trap paradox & 05 & open & 4 \\
13 & Filter-class dependence of trap economics & 05 & open & 4 \\
14 & DP-noised covariates through the Hill link & 06 & open & 5 \\
15 & Endogenous share drift & 06 & open & 5 \\
16 & Optimal release granularity G & 06 & open & 5 \\
17 & Analytic transport de-biasing & 06 & open & 5 \\
18 & Rotation theory of refresh decorrelation & 07 & open & 6 \\
19 & Posterior mean vs MAP flip rates & 07 & open & 6 \\
20 & Optimal hysteresis deadband h* & 07 & resolved by \dive{08} & 6 \\
21 & Multichannel radial deadband & 08 & resolved by \dive{09} & 7 \\
22 & Estimating (q, K) in the field & 08 & partly advanced by \dive{09} & 7 \\
23 & Robust filtering under heavy-tailed refresh noise & 08 & open & 7 \\
24 & Trust as a stock & 08 & open & 7 \\
25 & General-n tiered reallocation & 09 & open & 7 \\
26 & Adaptive-M stability & 09 & open & 7 \\
27 & The discrete-monitoring cost constant & 09 & open & 7 \\
28 & Cadence under measured refresh-error correlation & 09 & open & 7 \\
29 & Sharp segment-level $\Lambda$-bounds & 10 & open & 8 \\
30 & Ratio-balancing specification test & 10 & open & 8 \\
31 & Fixed-policy estimands under budget-coupled arms and frequency & 10 & open & 8 \\
32 & n-channel leakage matrix under spline time effects & 11 & open & 9 \\
33 & Drifting misallocation $\eta$\_gt & 11 & open & 9 \\
34 & The Hill bracket for imputed channels & 11 & open & 9 \\
35 & Auxiliary-share instruments & 11 & open & 9 \\
36 & Audit law under scale (basis) errors & 11 & open & 9 \\
37 & Boundary distribution of the profile-LR on the completely-monotone cone & 12 & open & 10 \\
38 & The (G, H) frontier of the long-run identified set & 12 & open & 10 \\
39 & Shape families for experiment-grounded de-confounding & 12 & open & 10 \\
40 & Onset-timing as a second surrogacy index & 12 & open & 10 \\
41 & Interval-valued posteriors in the deadband & 12 & open & 10 \\
42 & Proof of the layered-design theorem & 13 & resolved by \dive{14} & 11 \\
43 & Single-index identified set for attribution & 13 & open & 11 \\
44 & Cross-journey table for paths and frequency & 13 & open & 11 \\
45 & Attribution-induced targeting equilibrium & 13 & open & 11 \\
46 & Test \& Roll for attribution & 14 & open & 11 \\
47 & The df = 1 theorem & 14 & open & 11 \\
48 & Compliance-aware cell algebra & 14 & open & 11 \\
49 & Adaptive cap-raising & 14 & open & 11 \\
50 & Paired (odd-part) allocation & 14 & open & 11 \\
51 & Heavy-tailed certification & 15 & open & 12 \\
52 & Temporal-split contamination power law & 15 & open & 12 \\
53 & Multichannel counterfactual bar & 15 & open & 12 \\
54 & Certifiable bar for an existing library & 15 & open & 12 \\
55 & Calibrated-$\delta$ policy-agnostic CS & 16 & open & 13 \\
56 & Stall detection statistic & 16 & partly resolved by \dive{19} & 13 \\
57 & Multichannel stall geometry & 16 & resolved by \dive{19} & 13 \\
58 & Stall under adstock and spend-demand confounding & 16 & resolved by \dive{19} & 13 \\
59 & Bounded CEP-link memory state & 17 & open & 14 \\
60 & The survey-refresh $\delta$ pilot and its exclusion & 17 & open & 14 \\
61 & Panel-measured kernel as MMM prior & 17 & open & 14 \\
62 & Attention as an instrumented dose & 17 & open & 14 \\
63 & The $\kappa$ library & 18 & open & 15 \\
64 & Three-margin bound & 18 & open & 15 \\
65 & User-level Trimmed Match & 18 & open & 15 \\
66 & Exceedance-grid guard & 18 & open & 15 \\
67 & Regret-optimal cap under an informative prior & 18 & open & 15 \\
68 & Log-only confounding detector & 19 & open & 13 \\
69 & In-loop exploration under error-correlated wobble & 19 & open & 13 \\
70 & The budget-direction escape & 19 & priced by \dive{21} & 13 \\
71 & Switch cones under synergy & 19 & open & 13 \\
72 & Exact echo false-alarm law & 19 & open & 13 \\
73 & The n-channel spectral comb & 20 & resolved by \dive{21} & 1 \\
74 & Estimating Psi and S\_eps(omega) in the field & 20 & open & 1 \\
75 & Price the prior & 20 & resolved by \dive{22} & 1 \\
76 & Drift-capped identification budget & 20 & open & 1 \\
77 & Planner-feasible probes & 20 & open & 1 \\
78 & The window-contrast term in curvature information & 20 & open & 1 \\
79 & Intermodulation under synergy & 21 & open & 1 \\
80 & Bin assignment as an optimisation & 21 & open & 1 \\
81 & Audit a real flighting calendar & 21 & open & 1 \\
82 & The estimability audit & 21 & discharged for dives 20--22 by \dive{22} & 1 \\
83 & Closed form for the full-model optimal phase & 21 & open & 1 \\
84 & Delayed-peak adstock and the DC null & 21 & open & 1 \\
85 & The response-curve estimand for prior pricing & 22 & open & 1 \\
86 & The multichannel shape-prior ledger & 22 & open & 1 \\
87 & The EB widening rule as a shipped diagnostic & 22 & open & 1 \\
88 & The nonlinear ladder gap & 22 & open & 1 \\
\end{longtable}

\normalsize

\chapter{The program's protocol}
\label{app:protocol}

The rules below are those recorded in the program's index file; they are reproduced so that the epistemic
labels used throughout this compendium can be read against the standard each dive was held to.

\begin{enumerate}[leftmargin=2em]
\item Work the queue in order, one item per run. Sequence numbers are permanent and zero-padded.
\item Every fourth run is an EXTEND run: instead of a new item, revisit the prior dive with the most promising
``next steps'' section and push it further (deeper derivation, better simulation, stronger counterexample,
more complete design). Extensions get their own sequence number with the same slug and the suffix \texttt{\_ext}.
\item Depth standard (full protocol, as of run 03): formalise mathematically with explicit assumptions; run a
parallel literature subagent at dive start; a minimum of three construct--attack--refine rounds, iterating
until two consecutive rounds produce no material revision; at least eight numerical experiments including at
least two robustness sweeps and at least one negative or power control; uncertainty estimates on every
headline number and a devil's-advocate paragraph against every headline claim; a clean-room verification
subagent re-derives at least one core result before write-up; established results are separated from
speculation with clear labels. Write-ups include the attack-and-refine history.
\item Each dive ends with ``Next steps for a future session''---the hook extensions build on.
\item If a dive discovers a new open problem, append it to the queue backlog with a one-line rationale.
\end{enumerate}

Two consequences for the reader. Dives 01 and 02 predate the full protocol (rule 3 was adopted ``as of run
03''): they have no attack-and-refine history and no clean-room verification, which the reviewer notes of
\cref{ch:identification,ch:transport} take into account. And the ``clean-room'' verifier is a second
automated agent given only the claim and the model, not an independent human; its reproductions to machine
precision establish that a derivation is internally consistent under the stated model, not that the model
describes any market.

\chapter{File inventory}
\label{app:files}

Everything in this compendium derives from the following files, as they stood on the user's disk at
2026-09-07 (timestamps in \cref{tab:provenance}). Paths are relative to
\texttt{marketing\_measurement/frontier\_study/}.

\begingroup\raggedright
\begin{itemize}[leftmargin=1.2em,itemsep=1pt]
\item Framing documents: \texttt{README.md}, \texttt{01\_state\_of\_the\_field.md}, \texttt{02\_open\_questions.md}, \texttt{03\_mmm\_adoption\_barriers.md}.
\item Bibliographies: \texttt{resources/math\_stats.md}, \texttt{resources/engineering\_scale.md}, \texttt{resources/human\_wisdom.md}, \texttt{resources/mmm\_adoption\_barriers.md}, \texttt{resources/repos\_and\_tools.md}, \texttt{resources/books\_and\_institutes.md}.
\item Program index: \texttt{deep\_dives/00\_INDEX.md} (queue, backlog, log).
\item Dive reports: \texttt{deep\_dives/01\_mmm\_identified\_set.md} through \texttt{deep\_dives/22\_price\_of\_the\_prior.md} (twenty-two files; see \cref{tab:provenance}).
\item Code and results: \texttt{deep\_dives/code/NN\_<slug>.py} for each dive, with result files \texttt{17\_results\_quick.json}, \texttt{18\_results\_quick.json}, \texttt{19\_results\_all.json}, \texttt{22\_results.json}, the \texttt{20\_*.json} experiment outputs, and the scratch directories \texttt{\_scratch19/} and \texttt{\_scratch22/}. The code was consulted to disambiguate formulas and numbers; nothing was re-run for this compendium.
\end{itemize}
\endgroup
